\documentclass[preview]{standalone}

\usepackage{amsmath}
\usepackage{amssymb}
\usepackage{stellar}
\usepackage{definitions}

\begin{document}

\id{la-misurazione-in-psicologia}
\genpage

\section{La misurazione in psicologia}

\begin{snippetdefinition}{misura-psicologia-definition}{Misura}
    Per \emph{misura} si intende l'assegnazione di un numero o di un'etichetta
    alle caratteristiche degli oggetti, in base a regole definite e alla
    grandezza assunta come campione.
\end{snippetdefinition}

\begin{snippet}{scale-misurazione-psicologia}
    Nel 1951 Stevens distingue quattro tipi di scale di misurazione:
    \begin{itemize}
        \item \emph{scala nominale}, che classifica sulla base della presenza o
        assenza di una qualità;
        \item \emph{scala ordinale}, che permette graduatorie ma non informa
        sulla distanza tra gli eventi;
        \item \emph{scala a intervalli}, che mantiene costante l'intervallo tra
        posizioni successive;
        \item \emph{scala a rapporti}, che possiede le proprietà della scala a
        intervalli e include uno zero assoluto.
    \end{itemize}
\end{snippet}

\begin{snippetdefinition}{attendibilita-definition}{Attendibilità}
    L'\emph{attendibilità} si riferisce alla coerenza o stabilità dei dati
    comportamentali ottenuti attraverso test psicologici. Un risultato è
    attendibile se, a partire da condizioni simili, si ripete in tempi
    differenti.
\end{snippetdefinition}

\begin{snippetdefinition}{validita-definition}{Validità}
    La \emph{validità} indica il livello di precisione con cui un test o uno
    strumento di ricerca misura la variabile psicologica che intende misurare.
\end{snippetdefinition}

\begin{snippet}{tipi-validita}
    Si distinguono quattro tipi principali di validità:
    \begin{itemize}
        \item \emph{validità interna}, cioè il grado in cui le conclusioni
        permettono di ritenere esistente una relazione causa-effetto;
        \item \emph{validità esterna}, cioè il grado di generalizzazione dei
        risultati ad altri soggetti e situazioni;
        \item \emph{validità di costrutto}, cioè l'accuratezza con cui i
        risultati sono conformi alla teoria;
        \item \emph{validità statistica}, cioè la probabilità che la relazione
        tra variabili sia significativa e non casuale.
    \end{itemize}
\end{snippet}

\begin{snippetdefinition}{misure-self-report-definition}{Misure self-report}
    Le \emph{misure self-report} sono strumenti di autovalutazione, come
    questionari e interviste, attraverso cui i partecipanti forniscono
    informazioni su esperienze, stati interni, comportamenti o abitudini.
\end{snippetdefinition}

\begin{snippetdefinition}{questionario-definition}{Questionario}
    Un \emph{questionario} consiste in una serie di domande scritte che possono
    riguardare fatti, abitudini, stati affettivi, comportamenti passati o
    presenti. Può includere domande aperte, che permettono al soggetto di
    rispondere liberamente.
\end{snippetdefinition}

\begin{snippetdefinition}{intervista-interattiva-definition}{Intervista interattiva}
    Un'\emph{intervista interattiva} è un colloquio tra ricercatore e individuo
    finalizzato alla raccolta di informazioni dettagliate.
\end{snippetdefinition}

\begin{snippetnote}{limiti-self-report-note}{Limiti degli strumenti self-report}
    L'attendibilità e la validità degli strumenti self-report non sono
    garantite: i partecipanti possono fraintendere le domande, non ricordare
    chiaramente le proprie esperienze o rispondere in modo influenzato dalla
    desiderabilità sociale.
\end{snippetnote}

\includesnpt[width=65\%,src=/snippet/static-psychology/anxiety-self-report-scale.jpg]{centered-img}

\begin{snippetdefinition}{desiderabilita-sociale-definition}{Desiderabilità sociale}
    La \emph{desiderabilità sociale} è la tendenza a fornire risposte false o
    ingannevoli per creare un'immagine favorevole, o sfavorevole, di sé.
\end{snippetdefinition}

\end{document}
